\documentclass[12pt]{article}
\usepackage[utf8]{inputenc}
\usepackage[T1]{fontenc}
\usepackage{hyperref}
\title{Mubami Language Resources for K-12 Curriculum Development}
\author{Compiled by Claude Sonnet 5 (claude-sonnet-5) for Andr\'as Kornai}
\date{}
\begin{document}
\maketitle

\section{Introduction}
Mubami (ISO 639-3 tsx; Glottocode muba1238) is a Papuan language of the Western Province of Papua New Guinea, with about 1,700 speakers. Other names found in search results are Dausame, Tao-Suamato, Tao-Suame and Ta. Speakers live in villages along the Guavi and Aramia rivers. The search found very little material, and nothing designed for schools.

\section{Primary Educational Resources}
No primers, textbooks or school programmes for Mubami were found in search results.

\section{Linguistic and Reference Materials}
\begin{itemize}
\item Wikipedia overview of the Mubami language, which mentions word lists by Z'graggen (1975) and Reesink (1976): \url{https://en.wikipedia.org/wiki/Mubami_language}
\item ASJP (Automated Similarity Judgment Program) word list entry: \url{https://asjp.clld.org/languages/MUBAMI}
\end{itemize}

\section{Religious and Community Resources}
\begin{itemize}
\item Joshua Project profile of the Tao-Suamato people. It reports that a Bible translation has been requested but was still pending in that source: \url{https://joshuaproject.net/people_groups/15263}
\end{itemize}

\section{Digital and Audio Resources}
\begin{itemize}
\item Global Recordings Network page for the language. It lists a ``Dausuami'' variety and I have not checked what audio it holds: \url{https://globalrecordings.net/en/language/tsx}
\end{itemize}

\section{Curriculum Development Context}
Mubami is used by all age groups, but no written teaching materials were found. A curriculum would have to begin with new documentation, such as an orthography, a vocabulary and recorded texts, made together with the community.

\end{document}